The preceding sections established the current RECTOR stack and benchmark evidence.  This section turns that evidence into a concrete roadmap: what should be built next, what should be evaluated next, and which artifacts would make the system more useful in broader research and deployment-oriented settings.

%-------------------------------------------------------------------------------
\subsection{Deployment Priorities}
\label{subsec:limitations}
%-------------------------------------------------------------------------------

The most immediate next steps fall into five engineering tracks.

\subsubsection{Expand Rule Coverage}
The first priority is to close the gap between the training-time proxy layer and the full 28-rule audit layer.  Four rules currently remain audit-only because they require intent modeling, counterfactual reasoning, or multi-agent negotiation.  Extending proxy coverage, or adding conservative abstention mechanisms for those cases, would make the selection layer more complete and reduce the gap between optimization and audit.

\subsubsection{Improve Applicability and Generator Quality}
The second priority is to improve the quality of the trajectories being ranked and the quality of the rule-activation signal used to rank them.  The conservative Safety/Legal activation study, oracle-applicability cross-check, and selector ablations already identify the main leverage points: better recall on high-priority rules, stronger calibration for moderate-support rules, and generators that place more candidate mass inside the compliant region of trajectory space.

\subsubsection{Extend the Planning Horizon}
The third priority is to move beyond a single 5-second open-loop decision.  Strategic positioning, longer merges, and lane-preparation behaviors require either hierarchical planning or receding-horizon replanning.  Appendix~\ref{app:receding_horizon} outlines one such integration path with a pretrained M2I DenseTNT stack and provides a practical starting point for longer-horizon experimentation.

\subsubsection{Run Reactive Closed-Loop Evaluation}
The fourth priority is to evaluate RECTOR with a reactive closed-loop stack rather than only with open-loop replay and bridge-validation experiments.  The current Waymax integration already exercises the end-to-end execution pathway, so the next step is to replace the mock generator with the trained selector stack and pair it with reactive traffic participants.

\subsubsection{Broaden the Rule Catalog and Test Matrix}
The fifth priority is broader coverage: more jurisdictions, more map regimes, more weather and lighting conditions, and stronger perturbation testing for signal-state and lane-boundary uncertainty.  The verification results already identify the most sensitive inputs, which makes this expansion plan concrete rather than generic.

%-------------------------------------------------------------------------------
\subsection{Future Work}
\label{subsec:future_work}
%-------------------------------------------------------------------------------

It is useful mainly as an implementation sketch for longer-horizon and planning-loop work.

The roadmap above suggests a natural progression of future work, which we organize into near-term priorities and longer-term research directions.

\subsubsection{Near-Term Priorities}

The most pressing improvements are to expand proxy coverage for rules requiring intent modeling, extend the planning horizon via hierarchical or receding-horizon approaches, evaluate cross-region generalization, integrate prerequisite validity into applicability gating, and run full-scale reactive closed-loop studies with the trained RECTOR stack.

\subsubsection{Longer-Term Research Directions}

Beyond these immediate priorities, five research directions emerge from the current work.  \emph{Interaction validity}: extend the validated Waymax bridge into a full reactive-agent study with the trained RECTOR stack.  \emph{Rule learning}: learn intent and semantic predictors for semantics-heavy constraints and convert them into conservative applicability and severity bounds.  \emph{Adaptive parameterization}: support context-dependent rule thresholds while preserving deterministic priority ordering.  \emph{Jurisdictional parameterization}: develop versioned, auditable rule catalogs with regression suites for multi-region deployment.  \emph{Certification artifacts}: generate active-rule logs, normalized severities, first-differing-tier reports, and prerequisite-validity proofs as structured evidence.

%-------------------------------------------------------------------------------
\subsection{Open Questions}
\label{subsec:open_questions}
%-------------------------------------------------------------------------------

Finally, we identify five open questions raised by the current work but not resolved by it.  First, are 28 rules sufficient for a target operational design domain, and how should exceptions or contrary-to-duty rules be formalized?  Second, what constitutes a formally justified resolution policy when no candidate satisfies the highest-priority tier?  Third, how much of the accuracy--compliance gap is intrinsic to the selection mechanism versus an artifact of current generators?  Fourth, what minimal closed-loop assumptions suffice for safety evidence?  And fifth, how can rule catalogs be updated with traceability and without full retraining?
